\section{Separated finite-state automata from regular expressions}

In this section, we show that finite-state automata constructed from regular
expressions using the algorithm described in \Cref{thm:thompson} are $k$-separated
for $k > 1$. We first prove that the graph $\mathcal{G}_{\mathcal{A}}$ corresponding
to such finite-state automaton $\mathcal{A} = \langle \Sigma, Q, S, F, \Delta
\rangle$ is planar. As a consequence, the algorithm described in
\Cref{thm:lipton_tarjan} can be applied to find a separator of size
$\mathcal{O}(\sqrt{|Q|})$ in $\mathcal{O}(|Q|)$ time. We then describe a specialized
algorithm that is guaranteed to find a separator of constant size in
$\mathcal{O}(|Q|)$ time. An example of such separator is presented in
\Cref{figure:separator_thompson_automaton}.

\begin{figure}[!ht]
  \centering
  \begin{tikzpicture}[on grid, auto, scale=0.35]
    \node[state, initial, fill=white, scale=0.375]
    (s1) at (-14, 0) {};
    \node[state, fill=white, scale=0.375]
    (s2) at (-10, 0) {};
    \node[state, fill=black, scale=0.375]
    (s3) at (-6, 2) {};
    \node[state, fill=white, scale=0.375]
    (s4) at (-6, -2) {};
    \node[state, fill=white, scale=0.375]
    (s5) at (-2, 2) {};
    \node[state, fill=white, scale=0.375]
    (s6) at (-2, -2) {};
    \node[state, fill=white, scale=0.375]
    (s7) at (2, 2) {};
    \node[state, fill=white, scale=0.375]
    (s8) at (2, -2) {};
    \node[state, fill=black, scale=0.375]
    (s9) at (6, 2) {};
    \node[state, fill=white, scale=0.375]
    (s10) at (6, -2) {};
    \node[state, fill=white, scale=0.375]
    (s11) at (10, 0) {};
    \node[state, accepting, fill=white, scale=0.375]
    (s12) at (14, 0) {};

    \draw [->]
    (s1) edge node{$\epsilon$} (s2)
    (s2) edge node{$\epsilon$} (s3)
    (s2) edge node{$\epsilon$} (s4)
    (s3) edge node{$0$} (s5)
    (s5) edge node{$\epsilon$} (s7)
    (s7) edge node{$1$} (s9)
    (s4) edge node{$1$} (s6)
    (s6) edge node{$\epsilon$} (s8)
    (s8) edge node{$0$} (s10)
    (s9) edge node{$\epsilon$} (s11)
    (s10) edge node{$\epsilon$} (s11)
    (s11) edge node{$\epsilon$} (s12)
    (s1) edge[bend right=40] node{$\epsilon$} (s12)
    (s11) edge[bend left=-80] node[above]{$\epsilon$} (s2);
  \end{tikzpicture}
  \caption{Separator (colored in black) in a finite-state automaton derived from the
  regular expression $(((0\;.\;1) + (1\;.\;0))*)$ by applying the construction
  described in \Cref{thm:thompson}.}
  \label{figure:separator_thompson_automaton}
\end{figure}

\begin{proposition}
  Let $\mathcal{E}$ be a regular expression over an alphabet $\Sigma$ and let
  $\mathcal{A}$ be the finite-state automaton constructed from $\mathcal{E}$ by
  applying the algorithm described in $\Cref{thm:thompson}$. Then
  $\mathcal{G}_\mathcal{A}$ is planar.
\end{proposition}

\begin{proof}
  We show by induction on the construction of the finite-state automaton
  $\mathcal{A}$ that its corresponding graph $\mathcal{G}_{\mathcal{A}}$ can be
  embedded in a square so that the following properties hold.
  \begin{enumerate}
    \item Only the vertex corresponding to the single start state of $\mathcal{A}$ is
          mapped to a point that lies on the left side of the square,
    \item only the vertex corresponding to the single final state of $\mathcal{A}$ is
          mapped to a point that lies on the right side of the square,
    \item the vertices corresponding to all other states of $\mathcal{A}$ are mapped
          to points that lie properly inside the square,
    \item edges that correspond to transitions of $\mathcal{A}$ are mapped to simple
          curves that lie properly inside the square,
    \item curves intersect only at points corresponding to vertices of
          $\mathcal{G}_{\mathcal{A}}$.
  \end{enumerate}

  \begin{case}[$\mathcal{E} = \emptyset$]
    Then $\mathcal{A}$ is a Type 1 finite-state automaton. The embedding of the
    graph $\mathcal{G}_{\mathcal{A}}$ in a square is presented in
    $\Cref{fig:f_thompson_1}$. There are two vertices and no edges in
    $\mathcal{G}_{\mathcal{A}}$. Properties 1 and 2 hold for
    $\mathcal{G}_{\mathcal{A}}$ by construction. Properties 3, 4 and 5 hold for
    $\mathcal{G}_{\mathcal{A}}$ trivially.

    \begin{figure}[!ht]
      \centering
      \begin{tikzpicture}[on grid, auto, scale=0.5]
        \draw[dotted] (-4, -4) rectangle (4, 4);

        \node[state, fill=white, scale=0.75] (s) at
        (-4, 0) {$s$};
        \node[state, fill=white, scale=0.75] (f)
        at (4, 0) {$f$};
      \end{tikzpicture}
      \caption{Embedding of the graph corresponding to a Type 1 finite-state
      automaton in a square.}
      \label{fig:f_thompson_1}
    \end{figure}
  \end{case}

  \begin{case}[$\mathcal{E} = \sigma$ where $\sigma \in \Sigma_\epsilon$]
    Then $\mathcal{A}$ is a Type 2 finite-state automaton. The embedding of the
    graph $\mathcal{G}_{\mathcal{A}}$ in a square is presented in
    $\Cref{fig:f_thompson_2}$. There are two vertices and a single edge in
    $\mathcal{G}_{\mathcal{A}}$. Properties 1, 2, 4 and 5 hold for
    $\mathcal{G}_{\mathcal{A}}$ by construction. Property 3 holds for
    $\mathcal{G}_{\mathcal{A}}$ trivially.

    \begin{figure}[!ht]
      \centering
      \begin{tikzpicture}[on grid, auto, scale=0.5]
        \draw[dotted] (-4, -4) rectangle (4, 4);
        \node[state, fill=white, scale=0.75] (s) at
        (-4, 0) {$s$};
        \node[state, fill=white, scale=0.75] (f)
        at (4, 0) {$f$};
        \path (s) edge node {} (f);
      \end{tikzpicture}
      \caption{Embedding of the graph corresponding to a Type 2 finite-state
      automaton in a square.}
      \label{fig:f_thompson_2}
    \end{figure}
  \end{case}

  \begin{case}[$\mathcal{E} = \mathcal{E}'.\mathcal{E}''$ where $\mathcal{E}'$ and
    $\mathcal{E}''$ are regular expressions]
    Then $\mathcal{A}$ is a Type 3 finite-state automaton. The embedding of the graph
    $\mathcal{G}_{\mathcal{A}}$ in a square is presented in
    $\Cref{fig:f_thompson_3}$. Let $\mathcal{A}'$ and $\mathcal{A}''$ be the
    finite-state automata constructed from $\mathcal{E}'$ and $\mathcal{E}''$
    respectively. Also, let properties 1 through 5 hold for
    $\mathcal{G}_\mathcal{A'}$ and $\mathcal{G}_\mathcal{A''}$. Since property 1
    holds for $\mathcal{G}_\mathcal{A'}$ and property 2 holds for
    $\mathcal{G}_\mathcal{A''}$ then, by construction, properties 1 and
    2 hold for $\mathcal{G}_\mathcal{A}$. Also, since property 3 holds
    for $\mathcal{G}_\mathcal{A'}$ and $\mathcal{G}_\mathcal{A''}$
    then, by construction, it also holds for $\mathcal{G}_\mathcal{A}$.
    Since property 4 holds for $\mathcal{G}_\mathcal{A'}$ and
    $\mathcal{G}_\mathcal{A''}$, and the edge connecting $f'$ and $s''$
    is mapped to a simple curve that lies properly inside the square
    then property 4 holds for $\mathcal{G}_\mathcal{A}$. Moreover, that
    edge is mapped outside the square embeddings of the graphs
    $\mathcal{G}_\mathcal{A'}$ and $\mathcal{G}_\mathcal{A''}$. Thus,
    since properties 4 and 5 hold for them, then property 5 also holds for
    $\mathcal{G}_\mathcal{A}$.

    \begin{figure}[!ht]
      \centering
      \begin{tikzpicture}[on grid, auto, scale=0.25]
        \draw[dotted] (-8, -8) rectangle (8, 8);

        \draw (-8, -3) rectangle (-2, 3);
        \node[state, fill=white, scale=0.75]
        (s') at (-8, 0) {$s'$};
        \node[state, fill=white, scale=0.75]
        (f') at (-2, 0) {$f'$};

        \draw (2, -3) rectangle (8, 3);
        \node[state, fill=white, scale=0.75]
        (s'') at (2, 0) {$s''$};
        \node[state, fill=white, scale=0.75]
        (f'') at (8, 0) {$f''$};

        \path (f') edge node {} (s'');

        \node[] at (-5, 0) {$\mathcal{G}_\mathcal{A'}$};
        \node[] at (5, 0) {$\mathcal{G}_\mathcal{A''}$};
      \end{tikzpicture}
      \caption{Embedding of the graph corresponding to a Type 3 finite-state
      automaton in a square.}
      \label{fig:f_thompson_3}
    \end{figure}
  \end{case}

  \begin{case}[$\mathcal{E} = \mathcal{E}'+\mathcal{E}''$ where $\mathcal{E}'$ and
    $\mathcal{E}''$ are regular expressions]
    Then $\mathcal{A}$ is a Type 4 finite-state automaton. The embedding of the
    graph $\mathcal{G}_{\mathcal{A}}$ in a square is presented in
    $\Cref{fig:f_thompson_4}$. Let $\mathcal{A}'$ and $\mathcal{A}''$ be the
    finite-state automata constructed from $\mathcal{E}'$ and $\mathcal{E}''$
    respectively. Also, let properties 1 through 5 hold for
    $\mathcal{G}_\mathcal{A'}$ and $\mathcal{G}_\mathcal{A''}$. Properties 1 and 2
    hold for $\mathcal{G}_{\mathcal{A}}$ by construction. Since property 3 holds
    for $\mathcal{G}_\mathcal{A'}$ and $\mathcal{G}_\mathcal{A''}$, and
    since no vertices are ever mapped to points that lie on the top or
    bottom sides of a square then, by construction, it also holds for
    $\mathcal{G}_\mathcal{A}$. Also, since properties 4 and 5 hold for
    $\mathcal{G}_\mathcal{A'}$ and $\mathcal{G}_\mathcal{A''}$, and since the edges
    that connect $s$ to $s'$ and $s''$, and $f$ to $f'$ and $f''$ are mapped to
    curves that lie properly inside the square then property 4 holds for
    $\mathcal{G}_\mathcal{A}$. Moreover, since these edges are mapped to simple
    curves that lie outside the square embeddings of $\mathcal{G}_\mathcal{A'}$ and
    $\mathcal{G}_\mathcal{A''}$, and since they intersect other curves only at points
    where vertices of the graph are mapped to, then property 5 holds for
    $\mathcal{G}_\mathcal{A}$.

    \begin{figure}[!ht]
      \centering
      \begin{tikzpicture}[on grid, auto, scale=0.25]
        \draw[dotted] (-8, -8) rectangle (8, 8);

        \draw (-3, 8) rectangle (3, 2);
        \node[state, fill=white, scale=0.75]
        (s') at (-3, 5) {$s'$};
        \node[state, fill=white, scale=0.75]
        (f') at (3, 5) {$f'$};

        \draw (-3, -8) rectangle (3, -2);
        \node[state, fill=white, scale=0.75]
        (s'') at (-3, -5) {$s''$};
        \node[state, fill=white, scale=0.75]
        (f'') at (3, -5) {$f''$};

        \node[state, fill=white, scale=0.75]
        (s) at (-8, 0) {$s$};
        \node[state, fill=white, scale=0.75]
        (f) at (8, 0) {$f$};

        \path (s) edge node {} (s')
        (s) edge node {} (s'')
        (f') edge node {} (f)
        (f'') edge node {} (f);

        \node[] at (0, 5) {$\mathcal{G}_\mathcal{A'}$};
        \node[] at (0, -5) {$\mathcal{G}_\mathcal{A''}$};
      \end{tikzpicture}
      \caption{Embedding of the graph corresponding to a Type 4 finite-state
      automaton in a square.}
      \label{fig:f_thompson_4}
    \end{figure}
  \end{case}

  \begin{case}[$\mathcal{E} = \mathcal{E}'*$ where $\mathcal{E}'$ is a regular
    expression]
    Then $\mathcal{A}$ is a Type 5 finite-state automaton. The embedding of the graph
    $\mathcal{G}_{\mathcal{A}}$ in a square is presented in
    $\Cref{fig:f_thompson_5}$. Let $\mathcal{A}'$ be the finite-state automaton
    constructed from $\mathcal{E}'$. Also, let properties 1 through 5 hold for
    $\mathcal{G}_\mathcal{A'}$. Properties 1 and 2 hold for
    $\mathcal{G}_{\mathcal{A}}$ by construction. Since property 3 holds for
    $\mathcal{G}_\mathcal{A'}$ then, by construction, it also holds for
    $\mathcal{G}_\mathcal{A}$. Also, since properties 4 and 5 hold for
    $\mathcal{G}_\mathcal{A'}$ and since the edges that connect $s$ to $s'$ and $f$,
    and $f'$ to $s'$ and $f$ are mapped to curves that lie properly inside the square
    then property 4 holds for $\mathcal{G}_\mathcal{A}$. Moreover, since these edges
    are mapped to simple curves that lie outside the square embedding of
    $\mathcal{G}_\mathcal{A'}$ and since they intersect other curves only at points
    where vertices of the graph are mapped to, then property 5 holds for
    $\mathcal{G}_\mathcal{A}$.

    \begin{figure}[!ht]
      \centering
      \begin{tikzpicture}[on grid, auto, scale=0.25]
        \draw[dotted] (-8, -8) rectangle (8, 8);

        \draw (-3, -3) rectangle (3, 3);
        \node[state, fill=white, scale=0.75]
        (s) at (-8, 0) {$s$};
        \node[state, fill=white, scale=0.75]
        (s') at (-3, 0) {$s'$};
        \node[state, fill=white, scale=0.75]
        (f') at (3, 0) {$f'$};
        \node[state, fill=white, scale=0.75]
        (f) at (8, 0) {$f$};

        \path (s) edge node {} (s')
        (f') edge node {} (f)
        (f') edge [bend right=110, looseness=2]
        node {} (s')
        (s) edge [bend right=50, looseness=1.5]
        node {} (f);

        \node[] at (0, 0) {$\mathcal{G}_\mathcal{A'}$};
      \end{tikzpicture}
      \caption{Embedding of the graph corresponding to a Type 5 finite-state
      automaton in a square.}
      \label{fig:f_thompson_5}
    \end{figure}
  \end{case}

  Properties 1 through 5 guarantee that vertices of the graph are mapped to
  points in the plane and that edges of the graph are mapped to simple curves
  in the plane so that these curves intersect only at points where vertices
  of the graph are mapped to. It follows that $\mathcal{G}_{\mathcal{A}}$ is
  planar.
\end{proof}

We now show a specialized algorithm that finds a separator of constant size in a
finite-state automaton constructed using the algorithm described in
\Cref{thm:thompson}. The main procedure is described in
\Cref{alg:thompson_separator}. We proceed with proving the correctness of the
algorithm and providing a bound for the size of the resulting separator.

\begin{algorithm}
  \caption{Finds a constant-size separator for the set of states of a finite-state
  automaton $\mathcal{A}$ constructed from a regular expression using the algorithm
  described in \Cref{thm:thompson}.}
  \label{alg:thompson_separator}

  \DontPrintSemicolon
  \KwIn{$\mathcal{A} = \langle \Sigma, Q, S, F, \Delta \rangle$}
  \KwOut{$\langle X, P_1, P_2 \rangle$}
  $\mathcal{A}_c \gets \mathcal{A}$\;
  $X \gets S_c \cup F_c$\;
  $R \gets \emptyset$\;

  \While{$|Q_c| - |X| > \frac{2}{3} |Q|$}{
    \Switch{$\mathcal{A}_c$}{
      \uCase{Type 3}{
        \uIf {$|Q_c'| \geq |Q_c''|$}{
          $X \gets S_c' \cup F_c'$\;
          $\mathcal{A}_c \gets \mathcal{A}_c'$\;
          $R \gets R \cup Q_c''$\;
        }
        \Else{
          $X \gets S_c'' \cup F_c''$\;
          $\mathcal{A}_c \gets \mathcal{A}_c''$\;
          $R \gets R \cup Q_c'$\;
        }
      }
      \uCase{Type 4}{
        \uIf {$|Q_c'| \geq |Q_c''|$}{
          $X \gets S_c' \cup F_c'$\;
          $\mathcal{A}_c \gets \mathcal{A}_c'$\;
          $R \gets R \cup Q_c'' \cup S_c \cup F_c$\;
        }
        \Else {
          $X \gets S_c'' \cup F_c''$\;
          $\mathcal{A}_c \gets \mathcal{A}_c''$\;
          $R \gets R \cup Q_c' \cup S_c \cup F_c$\;
        }
      }
      \uCase{Type 5}{
        $X \gets S_c' \cup F_c'$\;
        $\mathcal{A}_c \gets \mathcal{A}_c'$\;
        $R \gets R \cup S_c \cup F_c$\;
      }
    }
  }
  \Return $\langle X, Q_c \setminus X, R \rangle$
\end{algorithm}

\begin{proposition}\label{prop:thompson_sep_1}
  Let $\mathcal{A} = \langle \Sigma, Q, S, F, \Delta \rangle$ be a finite-state
  automaton constructed using the algorithm described in \Cref{thm:thompson}. Then,
  if $\mathcal{A}_c$ is Type 1 or Type 2 during the execution of
  \Cref{alg:thompson_separator} on input $\mathcal{A}$, the condition at line 4 is
  false.
\end{proposition}

\begin{proof}
  Since $\mathcal{A}$ is a finite-state automaton constructed using the algorithm
  described in \Cref{thm:thompson}, then $|Q| \geq 2$. If $\mathcal{A}_c$ is Type 1
  or Type 2, then $|Q_c| - |X| = 0$. Therefore, the condition at line 4 is false.
\end{proof}

\begin{proposition}\label{prop:thompson_sep_2}
  Let $\mathcal{A} = \langle \Sigma, Q, S, F, \Delta \rangle$ be a finite-state
  automaton constructed using the algorithm described in \Cref{thm:thompson}. Then
  \Cref{alg:thompson_separator} completes on input $\mathcal{A}$ in
  $\mathcal{O}(|Q|)$ time.
\end{proposition}

\begin{proof}
  $\mathcal{A}_c$ is set to $\mathcal{A}$ at line 1. If at any execution of
  the condition at line 4 the finite-state automaton $\mathcal{A}_c$
  is Type 1 or Type 2 then by \Cref{prop:thompson_sep_1} the loop
  terminates. Otherwise, $\mathcal{A}_c$ is updated depending on its type.
  Since $|Q_c|$ decreases after each iteration of the loop, eventually the
  condition at line 4 will become false and the loop will not be executed
  anymore. Therefore, \Cref{alg:thompson_separator} completes on
  input $\mathcal{A}$. Moreover, the loop will be executed at most $|Q|$
  times. Since lines 10, 14, 20, 24 and 29 represent unions of disjoint sets
  the total time required to compute $R$ is $\mathcal{O}(|Q|)$. The rest of
  the lines within the loop require a constant amount of time per iteration.
  Therefore, \Cref{alg:thompson_separator} completes in
  $\mathcal{O}(|Q|)$ time.
\end{proof}

\begin{proposition}\label{prop:thompson_sep_3}
  Let $\mathcal{A} = \langle \Sigma, Q, S, F, \Delta \rangle$ be a finite-state
  automaton constructed using the algorithm described in \Cref{thm:thompson}. Then $Q
  = Q_c \cup R$ and $Q_c \cap R = \emptyset$ during the execution of
  \Cref{alg:thompson_separator} on input $\mathcal{A}$.
\end{proposition}

\begin{proof}
  $Q_c$ is set to $Q$ at line 1 and $R$ is set to $\emptyset$ at line 3.
  Therefore, the first time the condition at line 4 is executed $Q = Q_c \cup
    R$ and $Q_c \cap R = \emptyset$. After each iteration of the loop $Q_c$ and
  $R$ are updated depending on the type of $\mathcal{A}_c$.

  \begin{case}[$\mathcal{A}_c$ is Type 3]
    Then either $|Q_c'| \geq |Q_c''|$ or $|Q_c'| < |Q_c''|$. Let us
    assume that $|Q_c'| \geq |Q_c''|$. Let $Q_b$ be the value of $Q_c$
    before line 9 is executed and let $R_b$ be the value of $R$ before
    line 10 is executed. Then after line 10 we have $Q_c = Q_c'$ and $R
      = R_b \cup Q_c''$. Also, $Q_b = Q_c' \cup Q_c''$ and $Q_c' \cap
      Q_c'' = \emptyset$. Since $Q = Q_b \cup R_b$ and $Q_b \cap R_b =
      \emptyset$ then $Q = Q_c \cup R$ and $Q_c \cap R = \emptyset$.
    Similar argument follows in case $|Q_c'| < |Q_c''|$.
  \end{case}

  \begin{case}[$\mathcal{A}_c$ is Type 4]
    Then either $|Q_c'| \geq |Q_c''|$ or $|Q_c'| < |Q_c''|$. Let us
    assume that $|Q_c'| \geq |Q_c''|$. Let $Q_b$ be the value of $Q_c$
    before line 19 is executed and let $R_b$ be the value of $R$ before
    line 20 is executed. Then after line 20 we have $Q_c = Q_c'$ and $R
      = R_b \cup Q_c'' \cup S_c \cup F_c$. Also, $Q_b = Q_c' \cup Q_c''
      \cup S_c \cup F_c$ and $Q_c' \cap (Q_c'' \cup S_c \cup F_c) =
      \emptyset$. Since $Q = Q_b \cup R_b$ and $Q_b \cap R_b = \emptyset$
    then $Q = Q_c \cup R$ and $Q_c \cap R = \emptyset$. Similar argument
    follows in case $|Q_c'| < |Q_c''|$.
  \end{case}

  \begin{case}[$\mathcal{A}_c$ is Type 5]
    Let $Q_b$ be the value of $Q_c$ before line 28 is executed and let
    $R_b$ be the value of $R$ before line 29 is executed. Then after
    line 29 we have $Q_c = Q_c'$ and $R = R_b \cup S_c \cup F_c$. Also,
    $Q_b = Q_c' \cup S_c \cup F_c$ and $Q_c' \cap (S_c \cup F_c) =
      \emptyset$. Since $Q = Q_b \cup R_b$ and $Q_b \cap R_b = \emptyset$
    then $Q = Q_c \cup R$ and $Q_c \cap R = \emptyset$.
  \end{case}

  \noindent
  Therefore, $Q = Q_c \cup R$ and $Q_c \cap R = \emptyset$ during the
  execution of \Cref{alg:thompson_separator} on input $\mathcal{A}$.
\end{proof}

\begin{proposition}\label{prop:thompson_sep_4}
  Let $\mathcal{A} = \langle \Sigma, Q, S, F, \Delta \rangle$ be a finite-state
  automaton constructed using the algorithm described in \Cref{thm:thompson}. Then
  $|X| = 2$ during the execution of \Cref{alg:thompson_separator} on input
  $\mathcal{A}$.
\end{proposition}

\begin{proof}
  Since $\mathcal{A}$ is a constructed using the algorithm described in
  \Cref{thm:thompson}, then $|S| = |F| = 1$. Thus, $|S_c| = |F_c| = 1$ and $|X| = 2$
  at line 1. Therefore, the first time the condition at line 4 is executed $|X| = 2$.
  Each time $X$ is updated within the loop it is set to the union of the start states
  and the final states of a sub-automaton of $\mathcal{A}$. Since every sub-automaton
  of $\mathcal{A}$ is also constructed using the algorithm described in
  \Cref{thm:thompson}, then $|X| = 2$ after each update. Therefore, $|X| = 2$ each
  time the condition at line 4 is executed and after the last iteration of the loop
  $|X| = 2$ at line 32.
\end{proof}

\begin{proposition}\label{prop:thompson_sep_5}
  Let $\mathcal{A} = \langle \Sigma, Q, S, F, \Delta \rangle$ be a finite-state
  automaton constructed using the algorithm described in \Cref{thm:thompson}. Then
  $|R| \leq \frac{2}{3} |Q|$ during the execution of \Cref{alg:thompson_separator} on
  input $\mathcal{A}$.
\end{proposition}

\begin{proof}
  At line 3 $R$ is set to $\emptyset$ and since $|Q| \geq 2$ it follows that
  $|R| \leq \frac{2}{3} |Q|$. After each iteration of the loop $|R|$ increases.
  Let $\mathcal{A}_f$ be the value of $\mathcal{A}_c$ and $R_f$ be the value
  of $R$ when the condition at line 4 is false. The loop will not be executed
  anymore so $|R|$ will not increase further. Let us assume that $|R_f| >
    \frac{2}{3} |Q|$. If $\mathcal{A}_f = \mathcal{A}$ then $|R_f| = 0$
  and since $|Q| \geq 2$ this contradicts the assumption. Thus, there exists
  an automaton $\mathcal{A}_b$ that is the value of $\mathcal{A}_c$ and a set
  $R_b$ that is the value of $R$ just before $\mathcal{A}_c$ is set to
  $\mathcal{A}_f$. Moreover, $\mathcal{A}_b$ cannot be Type 1 or Type 2.

  \begin{case}[$\mathcal{A}_b$ is Type 3]
    Then either $\mathcal{A}_f = \mathcal{A}_b'$ or $\mathcal{A}_f =
      \mathcal{A}_b''$. Let us assume that $\mathcal{A}_f =
      \mathcal{A}_b'$. Then $|Q_b'| \geq |Q_b''|$. Also, since the
    condition at line 4 is true for $Q_b$ it follows that $|Q_b| - 2 =
      |Q_b'| + |Q_b''| - 2 > \frac{2}{3} |Q|$. Therefore, $|Q_b'| >
      \frac{1}{3} |Q|$ and since $|Q_b'| + |Q_b''| + |R_b| = |Q|$ then
    $|R_b| + |Q_b''| = |R_f| < \frac{2}{3} |Q|$ which contradicts the
    assumption. Similar argument follows in case $\mathcal{A}_f =
      \mathcal{A}_b''$.
  \end{case}

  \begin{case}[$\mathcal{A}_b$ is Type 4]
    Then either $\mathcal{A}_f = \mathcal{A}_b'$ or $\mathcal{A}_f =
      \mathcal{A}_b''$. Let us assume that $\mathcal{A}_f =
      \mathcal{A}_b'$. Then $|Q_b'| \geq |Q_b''|$. Also, since the
    condition at line 4 is true for $Q_b$ it follows that $|Q_b| - 2 =
      |Q_b'| + |Q_b''| > \frac{2}{3} |Q|$. Therefore, $|Q_b'| >
      \frac{1}{3} |Q|$ and since $|Q_b'| + |Q_b''| + |R_b| + 2 = |Q|$ then
    $|R_b| + |Q_b''| + 2 = |R_f| < \frac{2}{3} |Q|$ which contradicts
    the assumption. Similar argument follows in case $\mathcal{A}_f =
      \mathcal{A}_b''$.
  \end{case}

  \begin{case}[$\mathcal{A}_b$ is Type 5]
    Then $\mathcal{A}_f = \mathcal{A}_b'$. Since the condition at line 4
    is true for $Q_b$ it follows that $|Q_b| - 2 = |Q_b'| > \frac{2}{3}
      |Q|$. Therefore, $|R_f| = |R_b| + 2 = |Q| - |Q_b'| < \frac{1}{3}$
    which contradicts the assumption.
  \end{case}

  \noindent
  Therefore, $|R| \leq \frac{2}{3} |Q|$ during the execution of
  \Cref{alg:thompson_separator} on input $\mathcal{A}$.
\end{proof}

\begin{proposition}\label{prop:thompson_sep_6}
  Let $\mathcal{A} = \langle \Sigma, Q, S, F, \Delta \rangle$ be a finite-state
  automaton constructed using the algorithm described in \Cref{thm:thompson}. Then
  $\Delta \cap (Q_c \setminus X \times \Sigma_\epsilon \times R) = \Delta \cap (R
  \times \Sigma_\epsilon \times Q_c \setminus X) = \emptyset$ during the execution of
  \Cref{alg:thompson_separator} on input $\mathcal{A}$.
\end{proposition}

\begin{proof}
  At line 3 $R$ is set to $\emptyset$. Therefore, before the first iteration
  of the loop at line 4 we have $\Delta \cap (Q_c \setminus X \times
    \Sigma_\epsilon \times R) = \Delta \cap (R \times \Sigma_\epsilon \times Q_c
    \setminus X) = \emptyset$. After each iteration of the loop $Q_c$, $X$ and
  $R$ are updated depending on the type of $\mathcal{A}_c$.

  \begin{case}[$\mathcal{A}_c$ is Type 3]
    Then either $|Q_c'| \geq |Q_c''|$ or $|Q_c'| < |Q_c''|$. Let us
    assume that $|Q_c'| \geq |Q_c''|$. Then $X = S_c' \cup F_c'$, $Q_c =
      Q_c'$ and $R = R \cup Q_c''$. The only transition in $\Delta \cap
      (Q_c \times \Sigma_\epsilon \times R)$ is $\langle f',
      \epsilon, s'' \rangle$ where $f' \in F_c', \; s'' \in S_c''$.
    However, $f' \not\in Q_c \setminus X$. Also, $\Delta \cap (R \times
      \Sigma_\epsilon \times Q_c) = \emptyset$. Therefore,
    $\Delta \cap (Q_c \setminus X \times \Sigma_\epsilon \times R) =
      \Delta \cap (R \times \Sigma_\epsilon \times Q_c \setminus X) =
      \emptyset$. Similar argument follows in case $|Q_c'| < |Q_c''|$.
  \end{case}

  \begin{case}[$\mathcal{A}_c$ is Type 4]
    Then either $|Q_c'| \geq |Q_c''|$ or $|Q_c'| < |Q_c''|$. Let us
    assume that $|Q_c'| \geq |Q_c''|$. Then $X = S_c' \cup F_c'$, $Q_c =
      Q_c'$ and $R = R \cup Q_c'' \cup S_c \cup F_c$. The only transition
    in $\Delta \cap (R \times \Sigma_\epsilon \times Q_c)$ is $\langle
      s, \epsilon, s' \rangle$ where $s \in S_c, \; s' \in S_c'$. The only
    transition in $\Delta \cap (Q_c \times \Sigma_\epsilon \times R)$ is
    $\langle f', \epsilon, f \rangle$ where $f' \in F_c', \; f \in F_c$.
    However, $s', f' \not\in Q_c \setminus X$. Therefore, $\Delta \cap
      (Q_c \setminus X \times \Sigma_\epsilon \times R) = \Delta \cap (R
      \times \Sigma_\epsilon \times Q_c \setminus X) = \emptyset$.
    Similar argument follows in case $|Q_c'| < |Q_c''|$.
  \end{case}

  \begin{case}[$\mathcal{A}_c$ is Type 5]
    Then $X = S_c' \cup F_c'$, $Q_c = Q_c'$ and $R = R \cup S_c \cup
      F_c$. The only transition in $\Delta \cap (R \times \Sigma_\epsilon
      \times Q_c)$ is $\langle s, \epsilon, s' \rangle$ where $s \in S_c,
      \; s' \in S_c'$. The only transition in $\Delta \cap (Q_c \times
      \Sigma_\epsilon \times R)$ is $\langle f', \epsilon, f \rangle$
    where $f' \in F_c', \; f \in F_c$. However, $s', f' \not\in Q_c
      \setminus X$. Therefore, $\Delta \cap (Q_c \setminus X \times
      \Sigma_\epsilon \times R) = \Delta \cap (R \times \Sigma_\epsilon
      \times Q_c \setminus X) = \emptyset$.
  \end{case}

  \noindent
  Therefore, $\Delta \cap (Q_c \setminus X \times \Sigma_\epsilon \times R) =
    \Delta \cap (R \times \Sigma_\epsilon \times Q_c \setminus X) = \emptyset$
  during the execution of \Cref{alg:thompson_separator} on input
  $\mathcal{A}$.
\end{proof}

\begin{theorem}
  Let $\mathcal{A} = \langle \Sigma, Q, S, F, \Delta \rangle$ be a finite-state
  automaton constructed using the algorithm described in \Cref{thm:thompson}.
  There exists an algorithm that, given $\mathcal{A}$, finds a separator $X$ and a
  partition $\{ X, Q_1, Q_2 \}_Q$ such that $|X| = 2$ and $|Q_i| \leq \frac{2}{3}
  |Q|$ for $i \in \{ 1, 2 \}$ in $\mathcal{O}(|Q|)$ time.
\end{theorem}

\begin{proof}
  Follows directly from \Cref{prop:thompson_sep_2,prop:thompson_sep_4,prop:thompson_sep_5,prop:thompson_sep_6}.
\end{proof}
